\section{A quartic with coefficient height one}\label{sec:heightone}

The principal new mechanism is a coefficient reduction that applies to any
finite quadratic system over the integers. It uses nonnegative quadratic
forms to enforce agreement between separate occurrences of a variable.
The separation is essential: algebraically identifying these occurrences
before collecting coefficients would destroy the height bound.

\subsection{The block identity}

Let $\sigma_1,\ldots,\sigma_k\in\{-1,1\}$. For two lists of real numbers
$m=(m_1,\ldots,m_k)$ and $w=(w_1,\ldots,w_k)$, define
\begin{align}
\mathcal B_\sigma(m,w)
  ={}&\sum_{i=1}^k(m_i^2+w_i^2-m_iw_i)\\[-2mm]
   &+\sum_{1\le i<j\le k}\sigma_i\sigma_j(m_im_j+w_iw_j).
   \label{eq:block}
\end{align}
The exact identity
\begin{equation}
 2\mathcal B_\sigma(m,w)=
 \left(\sum_i\sigma_i m_i\right)^2+
 \left(\sum_i\sigma_i w_i\right)^2+
 \sum_i(m_i-w_i)^2
 \label{eq:block-sos}
\end{equation}
proves both nonnegativity and the equivalence
\begin{equation}
 \mathcal B_\sigma(m,w)=0
 \quad\Longleftrightarrow\quad
 w_i=m_i\ (1\le i\le k),\qquad \sum_i\sigma_i m_i=0.
 \label{eq:block-zero}
\end{equation}
These statements hold over any ordered field, and consequently over
$\mathbb Z$ as well as $\mathbb R$. The factor $1/2$ in the proof is harmless:
the definition \eqref{eq:block} is already an integer polynomial.

Ordinary squaring of a residual produces coefficients $2$ in its mixed
terms. Formula \eqref{eq:block} distributes each diagonal contribution
between two matching constraints. After division by two, every displayed
coefficient is $1$ or $-1$. The remaining issue is to prevent distinct
displayed terms from becoming the same monomial.

\subsection{Separating occurrences and preserving degree four}

First write each source row as a list of signed monomial occurrences,
\begin{equation}
 R_j(x)=\sum_{i=1}^{k_j}\sigma_{ji}M_{ji}(x),
 \qquad \deg M_{ji}\le 2.
 \label{eq:unit-rows}
\end{equation}
A coefficient $c$ can be written as $|c|$ occurrences with sign
$\operatorname{sgn}(c)$. Section~\ref{sec:binary-normalization} gives a
more economical preliminary normalization for large coefficients.
Occurrences remain distinct even when their monomials agree in the
source row.

For each occurrence introduce two entirely fresh variables $a_{ji},b_{ji}$
and set $m_{ji}=a_{ji}b_{ji}$. Label the two factors as follows:
\begin{center}
\begin{tabular}{lll}
\toprule
Source monomial & Label of $a_{ji}$ & Label of $b_{ji}$\\
\midrule
$x_p x_q$ & $x_p$ & $x_q$\\
$x_p$ & $x_p$ & $1$\\
$1$ & $1$ & $1$\\
\bottomrule
\end{tabular}
\end{center}
Even for $x_p^2$, the two aliases are different variables with the same
label. No alias is used in two occurrences. Introduce one further fresh
variable $w_{ji}$ for each occurrence, and use the block
$\mathcal B_{\sigma_j}(m_j,w_j)$.

All $m_{ji}$ have degree exactly two, including padded linear and constant
occurrences. Thus the $m_{ji}m_{j\ell}$ terms have degree four,
the $m_{ji}w_{ji}$ terms have degree three, and the $w_{ji}w_{j\ell}$
terms have degree two. This deliberate separation of degrees is what
allows the agreement constraints below to be added without coefficient
collisions.

\subsection{Agreement cycles and the single unit anchor}

For a list $z_1,\ldots,z_d$ with $d\ge3$, put $z_{d+1}=z_1$ and define
\begin{equation}
 \mathcal C(z_1,\ldots,z_d)
   =\sum_{i=1}^d z_i^2-\sum_{i=1}^d z_i z_{i+1}
   =\frac12\sum_{i=1}^d(z_i-z_{i+1})^2.
 \label{eq:cycle}
\end{equation}
This polynomial is nonnegative and vanishes exactly when all its entries
agree. Its coefficients are $1$ and $-1$. The restriction $d\ge3$ is
necessary: at $d=2$ the two edge monomials coincide and give coefficient
$-2$.

Group all aliases carrying the same old-variable label $x_p$.
For a group of size at least three, add its cycle \eqref{eq:cycle}.
For a group of size two, add one fresh dummy alias with the same label
and use the resulting triangle. A singleton needs no agreement
constraint. An unused source variable can be omitted when only existence
of solutions matters; all source variables of the instantiated equation
occur.

List all aliases with label $1$ as $u_1,\ldots,u_d$. Add fresh unit
aliases if necessary so that $d\ge2$, and put
\begin{align}
 \mathcal U(u_1,\ldots,u_d)
   & =1+\sum_{i=1}^d u_i^2-u_1-u_d
       -\sum_{i=1}^{d-1}u_i u_{i+1}\label{eq:unit-anchor}\\
   & =\frac12\left[(1-u_1)^2+
       \sum_{i=1}^{d-1}(u_i-u_{i+1})^2+(u_d-1)^2\right].
       \nonumber
\end{align}
It vanishes exactly when every $u_i=1$. Only this block contributes a
constant term, namely $1$, or linear terms, namely $-u_1-u_d$.

\begin{definition}[The height-one polynomial]\label{def:heightone}
For the fixed occurrence lists \eqref{eq:unit-rows}, define
\begin{equation}
 \mathcal H=
 \sum_j\mathcal B_{\sigma_j}(m_j,w_j)
   +\sum_{\text{old-variable groups}}\mathcal C
   +\mathcal U.
 \label{eq:heightone}
\end{equation}
The second sum includes the prescribed triangles for groups of size two
and omits singletons. All sums are fixed and finite. All coordinates
range independently over $\mathbb Z$.
\end{definition}

\begin{theorem}[Coefficient and zero-set theorem]\label{thm:heightone}
For a finite quadratic integer system with at least one nonzero monomial
occurrence, the construction above
produces an ordinary polynomial $\mathcal H$ with the following properties:
\begin{enumerate}
\item Every nonzero coefficient of its fully collected expansion is
      $1$ or $-1$; its constant coefficient is $1$.
\item Its total degree is exactly four, and its degree in each individual
      variable is at most two.
\item It is nonnegative on every real tuple.
\item It has an integer zero if and only if all source rows have a
      simultaneous integer zero.
\item If every source variable occurs, projection to one representative
      of each old-variable alias group is a bijection on integer zero
      sets. Its inverse is explicitly polynomial.
\end{enumerate}
\end{theorem}

\begin{proof}
Nonnegativity follows from \eqref{eq:block-sos}, \eqref{eq:cycle}, and
\eqref{eq:unit-anchor}. Suppose that $\mathcal H=0$. Every block is
nonnegative, so every block is zero. The cycles make each old-variable
alias group constant; choose its common value as $x_p$. The anchor makes
all unit aliases equal to one. Consequently each $m_{ji}$ becomes its
source monomial $M_{ji}(x)$. Formula \eqref{eq:block-zero} then gives
$R_j(x)=0$ for every source row.

Conversely, given a simultaneous source zero $x$, set every old-variable
alias to its label's value, every unit alias to one, every agreement
dummy to its group's value, and $w_{ji}=M_{ji}(x)$. Every block vanishes.
All added values are uniquely forced on a zero, which proves the
bijection and its explicit inverse. Before coefficient normalization,
this inverse has degree at most two in the source variables.

For the coefficient claim, partition the collected monomials by degree
and by the variable families they contain. At degree four, the terms
$m_{ji}^2$ use exactly one private alias pair, each with exponent two;
the terms $m_{ji}m_{j\ell}$ use two distinct private pairs. Each pair
belongs to precisely one occurrence, so neither kind can collide with
another term. Across different rows all alias supports are disjoint.

At degree three, each $-m_{ji}w_{ji}$ contains its own private $w$ variable
and its private alias pair. These monomials are all different.
At degree two, the terms containing $w$ variables come only from the row
blocks. A square uses one $w$ index; a mixed term uses a pair from the
same row. They are distinct, and no $w$ is an alias. The remaining
quadratic terms are precisely the agreement cycles and unit anchor.
Their alias groups are disjoint. Every cycle has at least three distinct
vertices, so each undirected edge occurs once. The anchor path likewise
has no repeated edge. Its two linear terms are distinct because $d\ge2$.
The constant term occurs once. Thus every displayed term contributes to
a distinct monomial, proving height exactly one.

Any occurrence contributes $a_{ji}^2b_{ji}^2$ with coefficient one, so
degree four is attained. The same enumeration shows that no individual
exponent exceeds two.
\end{proof}

The polynomial is a rational sum of squares, through the identity for
$2\mathcal H$. An integer sum-of-squares representation of $\mathcal H$
itself is neither used nor asserted. The proof relies on the order of
$\mathbb Z$ or $\mathbb R$; reducing the single equation modulo a prime
does not preserve its zero-set interpretation.

\subsection{A small example: one addition equation}

For $x+y-z=0$, take three unit aliases $u_1,u_2,u_3$, put
$m_1=xu_1$, $m_2=yu_2$, $m_3=zu_3$, and introduce $w_1,w_2,w_3$.
Then
\begin{align*}
\mathcal H_{+}={}&\sum_{i=1}^3(m_i^2+w_i^2-m_iw_i)
 +m_1m_2-m_1m_3-m_2m_3\\
 &+w_1w_2-w_1w_3-w_2w_3
 +1+u_1^2+u_2^2+u_3^2-u_1-u_3-u_1u_2-u_2u_3.
\end{align*}
Its collected coefficients are all $\pm1$, its degree is four, and
$\mathcal H_+=0$ holds exactly when
\[
 x+y=z,\qquad (u_1,u_2,u_3)=(1,1,1),\qquad(w_1,w_2,w_3)=(x,y,z).
\]
This example shows why addition is compatible with height one. No
positivity restriction on $x,y,z$, or on any auxiliary variable, is
needed.

\subsection{Normalizing large coefficients}\label{sec:binary-normalization}

Expanding a coefficient of magnitude $c$ into $c$ occurrences is correct
but can be wasteful. The delivered compiler first constructs each needed
positive integer $c$ from a shared unit and smaller constants. For
$c=2d+\epsilon$ with $\epsilon\in\{0,1\}$, introduce a constant wire
$C_c$ and the row
\begin{equation}
 C_c-C_d-C_d-\epsilon=0.
 \label{eq:binary-constant}
\end{equation}
The repeated $C_d$ is retained as two signed occurrences. Constant wires
are ordered by their defining dependencies. Here $C_1$ denotes the
literal constant $1$ and requires no additional coordinate. The resulting
finite triangular system forces each
$C_c$ to equal its stated integer and has a unique extension of every
source tuple. Sharing the binary ancestors avoids repeating the same
constant construction.

For a coefficient multiplying a quadratic monomial $x_p x_q$, introduce
or reuse a scaled wire $S_{c,p}$ with
\begin{equation}
 S_{c,p}-C_c x_p=0,
 \label{eq:scaled-wire}
\end{equation}
and replace $c x_p x_q$ by $S_{c,p}x_q$. A linear term $c x_p$ may
be written directly as $C_c x_p$, and a constant $c$ as $C_c$.
Signs remain in $\sigma_{ji}$. All new rows are quadratic and use only
unit signed occurrences. Substituting the forced constants and scaled
wires restores each original source row exactly. This is a polynomial
identity, independently checked in the release, rather than a test on
sampled source zeros.

The binary normalization and the height-one construction compose. Every
source tuple has a uniquely specified extension to the new coordinates
at a zero. For the delivered normalization, its constant wires are
constant and its scaled wires are linear in the source variables; the
full inverse on zero sets therefore still has degree at most two.
In fact, if $L$ denotes this explicit extension, the stronger identity
\begin{equation}
 \mathcal H(L(x))=\sum_j R_j(x)^2
 \label{eq:forward-value}
\end{equation}
holds for every source tuple, not only for source zeros. The constant
and scaling rows vanish identically under $L$, each pair of row squares
in \eqref{eq:block-sos} contributes twice the original squared residual,
and all agreement and unit blocks vanish.

\subsection{Size accounting and the instantiated equation}

Let $L=\sum_j k_j$ be the number of signed occurrences after normalization.
The row blocks use $3L$ variables: two aliases and one $w$ per occurrence.
Let $d_2$ count old-variable labels appearing exactly twice among aliases,
and let $\delta_u\in\{0,1,2\}$ be the number of additional unit aliases
needed to make the unit group have size at least two. The exact count is
\begin{equation}
 N=3L+d_2+\delta_u.
 \label{eq:height-count}
\end{equation}
A block with $k$ occurrences contributes $k^2+2k$ distinct monomials:
$k$ quartic squares, $\binom{k}{2}$ quartic cross terms, $k$ cubic terms,
and $k+\binom{k}{2}$ quadratic terms. A $d$-vertex agreement cycle
contributes $2d$ monomials, and a $d$-unit anchor contributes $2d+2$.
These formulas are exact because the proof excluded collisions.

The instantiated source is the 519-integer, 321-row C222 system from
the incoming powered-saturation packet \cite{round14}, with its complete
quadratic residual list included unchanged. The artifact inventory and
measured counts appear in Section~\ref{sec:artifacts}. In particular,
the output contains every collected coefficient of $\mathcal H$; it
does not leave an input program, compiler choice, or coefficient oracle
to be supplied later.

The general construction changes the coefficient objective decisively:
height one is the least positive coefficient height of any nonzero
integer polynomial. This observation is only an optimum for that one
coordinate of the complexity problem. It makes no claim that the
resulting number of variables, support, or total description length is
minimal. The existing height-two construction \cite{heighttwo} has fewer
variables, and the existing native quartic has far fewer variables with
much larger coefficients. Those remain distinct tradeoffs.

\subsection{Logical consequence and precise novelty scope}

The construction is entirely finite. Its soundness and completeness use
identities of integer polynomials, finite dependency induction, and
nonnegativity of squares. These arguments are formalizable in ZFC.
Composing Theorem~\ref{thm:heightone} with the retained C222 calibration
therefore gives
\begin{equation}
 \ZFC\vdash
 \left[\forall z\in\mathbb Z^{\HeightOneVariables},\ \mathcal H(z)\ne0\right]
 \ \longleftrightarrow\ \Con(\ZFC).
 \label{eq:heightone-consistency}
\end{equation}
The height-one transformation does not introduce another unbounded
number-theoretic lemma. Its dependence on the inherited arithmetic
encoding is identified separately in Section~\ref{sec:calibration}.

More generally, quadratic circuit conversion followed by this theorem
reduces integer Diophantine solvability to multiquadratic quartics of
height one and constant coefficient one. This is a direct corollary of
the proved construction; it is not a claim to have newly established
the classical undecidability theorem. Searches of the frozen repository,
the current incoming packets, and selected literature did not locate
this particular block-and-cycle compiler. That bounded audit is not a
worldwide priority result. The concrete contribution is the proved
compiler, its independently checked fully instantiated consistency
polynomial, and its comparison with the supplied height-two endpoint.
